\section{Travaux connexes}

La détection automatique des fuites dans les réseaux de distribution d'eau a
été largement étudiée, avec des approches allant des techniques traditionnelles
de traitement du signal aux modèles modernes d'apprentissage automatique orientés
données. Les méthodes existantes peuvent être classées en trois grandes
catégories : approches basées sur la physique, méthodes classiques de machine
learning, et modèles fondés sur l'apprentissage profond.

\subsection{Méthodes traditionnelles et basées sur la physique}

Les premiers travaux en détection de fuites reposaient principalement sur la
modélisation hydraulique et les techniques de traitement du signal. Les méthodes
basées sur l'analyse des résidus de pression, la détection acoustique et
l'estimation du bilan de débit ont été largement utilisées en pratique. Ces
approches dépendent généralement d'une expertise métier et nécessitent des
modèles physiques précis du réseau d'eau. Bien qu'efficaces en environnement
contrôlé, elles rencontrent souvent des difficultés en conditions réelles en
raison du bruit, de l'incertitude et de la complexité des infrastructures
urbaines. De plus, leur dépendance à des seuils prédéfinis limite leur
adaptabilité aux conditions d'exploitation dynamiques.

\subsection{Approches de machine learning}

Avec la disponibilité croissante des données capteurs, des techniques de machine
learning ont été introduites pour la détection des fuites. Des modèles classiques
tels que les machines à vecteurs de support (SVM), les forêts aléatoires et le
gradient boosting ont été appliqués pour classifier les comportements normaux
et anormaux à partir de caractéristiques extraites. Ces approches améliorent le
passage à l'échelle par rapport aux systèmes à base de règles, mais nécessitent
souvent une ingénierie de caractéristiques importante et restent limitées dans
leur capacité à modéliser les dépendances temporelles inhérentes aux séries
temporelles.

Par ailleurs, les modèles classiques opèrent généralement sur des variables
agrégées ou construites manuellement, qui peuvent ne pas capturer la structure
séquentielle et l'évolution temporelle des événements de fuite. Cette limite
motive l'usage de modèles séquentiels capables d'apprendre directement à partir
de données temporelles brutes.

\subsection{Apprentissage profond pour la détection de fuites en séries temporelles}

Les avancées récentes en apprentissage profond ont permis une modélisation plus
efficace des séries temporelles multivariées. Les réseaux de neurones
convolutionnels (CNN) sont utilisés pour capturer les motifs temporels locaux,
tandis que les réseaux récurrents (RNN), en particulier les réseaux Long
Short-Term Memory (LSTM), sont capables de modéliser les dépendances de long
terme.

Plusieurs études ont exploré des architectures hybrides combinant CNN et LSTM
afin de tirer parti à la fois de l'extraction locale de caractéristiques et de
la modélisation temporelle. Ces modèles ont montré de meilleures performances
que les approches isolées, en particulier dans des environnements de surveillance
complexes où les signatures de fuite sont distribuées dans le temps.

Plus récemment, les mécanismes d'attention ont été introduits pour renforcer la
modélisation séquentielle en permettant aux modèles de se concentrer sur les pas
de temps les plus pertinents. L'attention a été appliquée avec succès dans
diverses applications de séries temporelles, améliorant à la fois les
performances et l'interprétabilité. Cependant, son application à la détection
de fuites d'eau reste encore limitée, et ses bénéfices n'ont pas été largement
validés par des comparaisons systématiques avec des modèles alternatifs.

\subsection{Limites des travaux existants}

Malgré ces avancées, plusieurs limites persistent. De nombreuses études ne
traitent pas correctement le déséquilibre de classes, qui constitue un défi
majeur en détection de fuites où les événements sont relativement rares. En
outre, l'évaluation est souvent limitée à des métriques centrées sur l'accuracy,
qui ne reflètent pas nécessairement les exigences opérationnelles réelles,
comme la minimisation des fausses alertes.

De plus, certaines approches ne garantissent pas une séparation stricte entre
données d'entraînement et de test dans un contexte temporel, ce qui peut mener
à des fuites de données et à des estimations de performance trop optimistes.
Une autre limite importante est l'absence de comparaisons de benchmark complètes,
ce qui rend difficile l'évaluation de l'efficacité réelle des modèles proposés
par rapport aux approches alternatives.

\subsection{Approche proposée}

Pour répondre à ces limites, cet article propose une architecture hybride
CNN-BiLSTM augmentée par un mécanisme d'auto-attention pour la détection de
fuites à partir de séries temporelles multivariées. Contrairement aux travaux
précédents, l'approche proposée intègre :

\begin{itemize}
    \item Une combinaison de CNN et BiLSTM pour capturer à la fois les motifs temporels locaux et de long terme,
    \item Un mécanisme d'auto-attention pour identifier les pas de temps les plus informatifs,
    \item Une stratégie de découpage temporel sans fuite de données,
    \item Un cadre d'évaluation adapté au déséquilibre de classes, utilisant l'AUC-PR et la calibration du seuil,
    \item Une comparaison de benchmark complète face à des modèles classiques et d'apprentissage profond.
\end{itemize}

Cette combinaison permet une détection robuste et fiable des fuites tout en
maintenant un compromis équilibré entre sensibilité de détection et taux de
fausses alertes, et fournit une évaluation plus complète de l'efficacité du
modèle en conditions réelles.